% !TeX root = ../papers/nc-p3.tex
\section{分岐中心に現れる純関係とその反対曲線}
\label{sec:branch-local-pure-pencils}
\begin{proposition}\label{prop:branch-local-pure-pencil}
曲面射影のない標準正則窓で分岐があれば、順・逆の少なくとも一つの中心の
三次元縮約像が純関係を含む。すなわち、順側なら順中心 $p$ の後続直線を消す二空間 $S$ に対して
$P=R_{i+1}\cap(S\otimes V_{i+2})$ に純関係があり、逆側なら語順を反転した条件が成り立つ。
\end{proposition}
\begin{proof}
全てのこの三空間が純関係を含まないと仮定する。
系\ref{cor:pure-free-branch-rank-nine}の鉛筆の議論は各三空間だけを使い、全中心が $L_3$ 型になる。
したがって各中心の後のHankel三次曲線、次中心の初回直線上への包含、
同位相二中心の制約は命題\ref{prop:pure-free-branch-cubic}のままである。
定理\ref{thm:branch-centers-opposite-phases}, \ref{thm:branch-cubic-contact}の
中心数・次数・交叉の証明はこれらと一般の次数公式しか使わない。
二直線列のない場合の最後の $(3,3)$ で、命題\ref{prop:branch-three-three-frame}が使う
第四の奇文字の切除による単射も保たれる。その核の元は第四の奇文字に支持される純関係であり、
その文字は初回・終回両直線を消すので、向きに従い順・逆の中心の三空間自体に入るからである。
従って定理\ref{thm:branch-pencil-resultant}が二直線列を強制する。

二直線列で $w(p,z)$ はどちらの端からも $L_3$ の非零縮約なので階数2である。
系\ref{cor:two-line-secant-roots}で使う $P\cap Q$ の上界1も、$P$ に純関係がないことだけを要する。
これで定理\ref{thm:two-line-double-tangency}, \ref{thm:pure-branch-degree-patterns}の全次数配置を得る。
最後の定理\ref{thm:final-pure-branch-frames}の四中心の場合に排除する一行支持の関係は
初回直線を消す中心の三空間に属する。三中心の場合は $U_1,U_2$ のうち第四行だけ、または
第四列だけに支持される関係を排除すればよい。位置 $1,2,3$ の第四文字はその初回・終回直線の両方を消すので、
それぞれ該当する順・逆中心の三空間に属する。従ってこれらも排除され、同じ二矛盾を得る。
\end{proof}

\begin{proposition}\label{prop:finite-pure-pencil-opposite}
全純関係スキームが有限で、$p$ が順中心であるとする。中央支持次元 $a$ は3か4である。
辺 $i+1$ の第一点 $x$ が初回直線 $L$ の外にあるとき、$P$ の三関係だけから得られる
後点候補は次の通りである。
\begin{enumerate}
\item $a=4$ なら $P$ の鉛筆は $L_e\oplus J_{3-e}$、$1\le e\le3$。
ここで $J$ は正方正則鉛筆で各固有値に対しJordanブロックは一つだけである。
一般の後点候補は $L$ からの射影の値に対応する有理正規 $e$ 次曲線の点であり、
例外の各方向ではその点を通る一直線に増える。
\item $a=3$ なら三行三列の正則鉛筆で、各固有値の幾何学的重複度は1である。
一般の後点候補は中央支持を消す定点、例外の各方向ではこの点を通る一直線である。
\end{enumerate}
三空間内の純関係スキームの長さは $a=4$ のとき $3-e$、$a=3$ のとき3である。
各位相に順中心は高々二つで、二つなら両方とも $a=4$、一方の列の第二次数は他方の $e$ である。
同じ位相で異なる中心の純関係は互いに素である。全て逆向きにも成り立つ。
\end{proposition}
\begin{proof}
$a=2$ のとき三次元の $2\times2$ 行列空間は正次元でSegreと交わる。
$S^*$ 方向 $[s:t]$ で $P$ を評価する $a\times3$ 鉛筆には定数の左右核がない
（両者とも実際の全支持を取ったからである）。
一般階数が3未満なら、各方向で少なくとも一つ非零の純関係が生じるため有限性に反する。
よって三行では全体が正方正則であり、四行ではKronecker標準形に左の特異ブロックが一つあり、
その最小指数は正である。各例外方向の核が二次元以上なら、同じ第一因子で純関係の直線ができる。
従って正方部の各固有値に一つのブロックしかない。
純関係に対しその第一因子を消す $[s:t]$ はスキーム上でも一意である
（純テンソルSegreの第一射影による）。その方向の三列の核線が純関係を復元する。
単元な二次小行列で核線を消去すれば、残る条件は三次小行列で、
局所Smith形によりその長さは正方部の対応する単因子の位数である。

後点は評価鉛筆の左核である。$L_e$ の左核は完全 $e$ 次系の点、三行正則なら
四文字内で使わなかった一方向が常に残る。特異方向ごとに正方部から一次元が加わる。
同位置の別の順中心の初回直線は $L$ と交わらず、そこからの $L$ を中心とする射影は同型なので、
その第二曲線は一般方向の後点軌跡に等しい。定点にはなれず、また異なる二つの初回移動が
同じ第二曲線を共有しない。よって中心数の上界と第二次数が従う。互いに交わらない初回直線の
消滅二空間の交叉は零なので三関係の空間も交わらない。
\end{proof}

\begin{proposition}\label{prop:excess-contact-two-line-pure}
二直線列 $p,L,M,z$ の始点にある別の動く曲線枝を $x(t)$ とする。
$p$ の点イデアルの引戻しの位数を $h$、その枝上の順評価の三次小行列式の
共通位数を $f$ とし、一般順階数は3とする。$f\ge h$ であり、$f>h$ なら
$R_{i+1}$ の $2\times 2$ 部分空間 $L^\perp\otimes M^\perp$ に独立な二関係がある。
従ってその辺の純関係スキームは少なくとも長さ2の部分スキームを含む。
逆向きにも同じである。
\end{proposition}
\begin{proof}
局所パラメータを $t$ とし、$R_i$ の六関係に対する $x(t)$ の評価の四行を考える。
$L$ の二座標で読む最初の二行はともに $t^h$ で割れる。
それらを $t^h$ で割って剰余体に落とした二行は独立である。
実際 $R_i$ の $L$ への制限は $\ker p\otimes H^0(L,\mathcal O(1))$ への同型で、
二行は $x(t)$ の最初の非定数方向と独立な二座標を対にしたものである。
他の二行は $p$ の順階数が2なのでそのままで剰余体上独立である。
三次小行列は必ず最初の行群を使うので $f\ge h$。$f>h$ なら最初の二行を割った
各々と後の二行で作った三次小行列が剰余体で消え、両二行の張る空間は等しい。

六次元の相補切断によりこれらの行を $R_{i+2}$ の二文字行列と同定する。
$w$ の $L\times M$ への制限は二項括弧と、$\ker p,\ker z$ 上の非退化な三次正方行列の積である。
従って割った二行を $M$ に制限すれば、$M$ の独立な二文字と、非定数方向により
$\ker z$ に定まる同一の非零文字との積となる。この二行が $p$ の評価の二行でもあるため、
$w(p)$ の末尾からの三次元縮約像を $M$ に制限すると次元1になる。
その二次元核は $L^\perp\otimes M^\perp$ に入る。この二次元行列空間の射影直線を
二次の行列式で切れば、全直線または長さ2の零因子になる。
\end{proof}
